
Soit @@M0@@ une pr\'esentation de Chow~: @@M1@@ est projectif sur @@M2@@\,; et @@M3@@ est surjectif. Si l'on admet (a), (b), (c) et (iii) pour @@M4@@, alors ils en r\'esultent pour @@M5@@ en vertu de 4.3 et 3.5. On peut donc supposer @@M6@@ projectif sur @@M7@@.

Soit @@M8@@ ample pour @@M9@@. Si pour un @@M10@@, @@M11@@ est ample, alors @@M12@@ est ample pour toute g\'en\'erisation @@M13@@ de @@M14@@ (EGA III 4.7.1). Gr\^ace \`a l'\'equivalence de 4.6 (i) et (iv), on a donc \'etabli (a).

Pour (b), (c) et 4.7 (iii), on peut supposer @@M15@@ int\`egre avec @@M16@@ point g\'en\'erique dans (b) et (c) (Pour (b) et (c), c'est clair, et pour (iii), cela r\'esulte de 3.4 (i) appliqu\'e au morphisme @@M17@@, o\`u les @@M18@@ sont les composantes irr\'eductibles de @@M19@@, munies de la structure r\'eduite induite).

Il existe un sch\'ema int\`egre @@M20@@ et un morphisme f.p.p.f. @@M21@@, tels que toute composante irr\'eductible @@M22@@ de @@M23@@, munie de la structure r\'eduite induite soit int\`egre sur @@M24@@ (XII 1.3). Pour prouver (b) et (iii), il suffit de le faire pour @@M25@@, (dans (b), par descente\,; dans (iii) par r\'ecurrence noeth\'erienne et 3.4 (i)).

Enfin, consid\'erons le morphisme surjectif canonique @@M26@@. Alors (b) et (iii) pour @@M27@@ r\'esultent de (c) et (iii) pour les @@M28@@, en vertu de 4.3, 3.5 et 3.7.

On s'est donc ramen\'e \`a prouver (c) et (iii), en supposant @@M29@@ int\`egre de point g\'en\'erique @@M30@@, et @@M31@@ projectif et int\`egre sur @@M32@@. Soit @@M33@@ un faisceau tr\`es ample pour @@M34@@. Alors pour tout @@M35@@, @@M36@@ et @@M37@@ ne d\'ependent pas de @@M38@@ (EGA III 7.9.4). Donc (c) r\'esulte de 4.6 (i) @@M39@@ (v), qui implique aussi que la famille des classes des faisceaux inversibles @@M40@@-\'equivalents \`a @@M41@@ sur les fibres de @@M42@@ n'a qu'un
